% ---------------------------------------------------------------------------
% Kapitel 8: Zusammenfassung und Ausblick
% Kurzfassung vom 04.10.2026, neu geschrieben aus der Langfassung
% (thesis/langfassung/kapitel08.tex, Confluence-Seite 58097676).
% ---------------------------------------------------------------------------

\chapter{Zusammenfassung und Ausblick}
\label{kap:zusammenfassung}

\section{Zusammenfassung}
\label{sec:gesamtzusammenfassung}

Die vorliegende Arbeit hat ein institutionelles E-Portfolio-Portal für die Hochschule Bremen entworfen und gebaut. Aus Forschungsliteratur und heuristischer Analyse einer Referenzplattform entstand ein Katalog aus 36 funktionalen und 24 nicht-funktionalen Anforderungen, die Konzeption und Realisierung prüften, ob er einer Umsetzung standhält, und vier Methoden mit elf Personen prüften das Ergebnis an vorab gesetzten Maßstäben. 30 funktionale Anforderungen sind erfüllt und sechs teilweise, von den nicht-funktionalen zehn erfüllt, elf teilweise und drei ungeprüft. Der SUS-Mittelwert liegt mit 87,3 um 19,3 Punkte über dem Zielwert von 68 (Lewis \& Sauro, 2018); ihm stehen 30 Usability-Probleme gegenüber, sieben davon schwerwiegend.

\subsection{Antwort auf die erste Teilfrage}
\label{subsec:antwort-tf1}

Die erste Teilfrage lautet: „Welche Anforderungen an Struktur, Inhalt und Workflow lassen sich aus der bestehenden Forschungsliteratur zu E-Portfolios ableiten?“ Zur Struktur liefert die Literatur zwei Anforderungen. Aus der institutionellen Kategorie und der offenen Leserschaft nach Lorenzo und Ittelson (2005) folgen der Zugang ohne Anmeldung und das Gewicht der externen Nutzenden; die Übertragung auf studentische Arbeiten ist ein Schritt der vorliegenden Arbeit. Aus den nach ihrem Bezugspunkt geordneten Komponenten von Mahasneh (2020) folgt die Trennung von Profilseite und Projekt. Zum Inhalt liefert dasselbe Modell das Raster, von dem das Portal sieben Komponenten abdeckt und drei veranstaltungsgebundene bewusst auslässt. Zum Workflow trägt die Bestimmung des institutionellen Portfolios als fokussierte Auswahl bei (Kahn, 2001, zitiert nach Lorenzo \& Ittelson, 2005), aus der die redaktionelle Freigabe folgt. Quer dazu begründet die Akzeptanzforschung (Davis, 1989; Shroff et al., 2011), dass Gebrauchstauglichkeit eine Frage des Entwurfs ist. Von den neun funktionalen Anforderungen mit Herkunft aus der E-Portfolio-Literatur sind acht erfüllt; der neunten (FA-16) fehlt allein der Studiengang der Person. Nicht aus der Literatur folgen die fünf Zustände des Freigabewegs, die Ableitung der Rolle aus der Hochschuladresse, die Ordnung der Hochschule und die kuratierte Startseite.

\subsection{Antwort auf die zweite Teilfrage}
\label{subsec:antwort-tf2-kurz}

Die zweite Teilfrage lautet: „Welche funktionalen und gestalterischen Merkmale fördern die Auffindbarkeit und Benutzerfreundlichkeit eines institutionellen Portfolio-Portals?“ Die Auffindbarkeit fördern sichtbar angebotene, verbindbare Ordnungsmerkmale, eine einzeln rücknehmbare Einschränkung und ein Zugang über die Person: Alle elf Personen schränkten drei Merkmale zugleich ein, im Median in 59 Sekunden, und alle fanden die Profilseite. Die Benutzerfreundlichkeit fördern kurze Wege für wiederkehrende Vorgänge und ein ohne Einweisung bedienbarer Editor. Nicht bestätigt hat sich, dass Trefferzahl und Zählmarken den Zustand der Suche ausreichend zeigen. Hinzu kommt eine Bedingung: Handlungen und Zustände müssen beschriftet sichtbar sein (H6), und dasselbe muss überall gleich heißen (H4). Gegenüber der Referenzplattform schließt das Portal die Lücke bei H7 im öffentlichen Bereich und bleibt bei H2 und H4 zurück.

\subsection{Antwort auf die dritte Teilfrage}
\label{subsec:antwort-tf3}

Die dritte Teilfrage lautet: „Wie lässt sich ein blockbasiertes Inhaltsmodell mit integriertem Redaktionsworkflow realisieren, das Studierenden eine flexible und medienreiche Projektpräsentation ermöglicht?“ Die Realisierung bestätigt vier Entscheidungen. Das Inhaltsmodell ist an das Dokumentmodell des Editors gebunden, sodass eine schmale Zuordnung zwischen Editor und Speicherung genügt. Der Inhalt eines Bausteins wird unverändert gespeichert, seine Art zusätzlich als abfragbares Merkmal geführt. Dateien liegen außerhalb des Dokuments und werden über einen eigenen Weg ausgeliefert, weil ihre Sichtbarkeit am Projektzustand hängt. Und die Bearbeitbarkeit ist an diesen Zustand gebunden und auf dem Server durchgesetzt, auch für Dateien; darin besteht die Integration des Workflows in das Inhaltsmodell. Für mehrere Bearbeitende verlangt das automatische Speichern ein Versionsfeld, das Konflikte meldet, und eine Bearbeitungssperre, die sie selten macht. Die Flexibilität tragen fünfzehn Bausteinarten nach dem Komponentenmodell von Mahasneh (2020). Dass Flexibilität nicht immer mit Benutzerfreundlichkeit einhergeht (Davis, 1989), zeigt der Test an einer Stelle (P4), und der Schrägstrichbefehl im Editor ist ein Fehler der Umsetzung (P30).

\subsection{Antwort auf die Forschungsfrage}
\label{subsec:antwort-ff}

Die Forschungsfrage lautet: „Welche funktionalen und gestalterischen Anforderungen muss ein hochschulisches E-Portfolio-Portal erfüllen, um studentische Projekte strukturiert zu präsentieren, auffindbar zu machen und von den beteiligten Nutzergruppen akzeptiert zu werden?“ Strukturiert präsentiert ein Portal, das zwischen Projekt und Person trennt, Projekte aus frei anordenbaren Bausteinen aufbaut und vor der Veröffentlichung redaktionell freigibt. Auffindbar macht es seinen Bestand, wenn es die Ordnung der Hochschule als verbindbare, einzeln rücknehmbare Merkmale anbietet, die wirkende Einschränkung bei den Treffern zeigt und einen Zugang über die Person eröffnet. Akzeptiert wird es, wenn sein Nutzen erkennbar ist und Zustände wie Handlungen beschriftet sichtbar sind. Bei Nützlichkeit und Benutzerfreundlichkeit liegen alle elf Befragten über dem Orientierungspunkt, und wie bei Davis (1989) hängt die Nutzungsabsicht stärker mit der Nützlichkeit zusammen (0,84 gegenüber 0,41). Übertragbar erscheinen der Zugang ohne Anmeldung, die Freigabe vor der Veröffentlichung, die Trennung von Projekt und Person und die sichtbare Beschriftung von Zustand und Handlung. An die Hochschule Bremen gebunden bleiben die Ausgestaltung des Freigabewegs, die Ableitung der Rolle, die kuratierte Startseite und die Ordnung aus fünf Fakultäten mit 75 Studiengängen.

\subsection{Grenzen der Arbeit}
\label{subsec:grenzen}

Die empirische Grundlage ist schmal: elf Personen in einer einzigen Testrunde, Redaktion, Administration und externe Nutzende mit je zwei Personen unter der Mindestzahl nach Nielsen (2000), keine Lehrenden, und in den Rollenblöcken erklärten die Karten den Ablauf. Die heuristische Evaluation stammt von zwei Begutachtenden, die das System selbst gebaut haben (Nielsen \& Molich, 1990). Beide Bögen sind nicht eigens validiert, der TAM-Bogen erfragt die Absicht nach einer Sitzung statt nach einem Semester, und die Teilnehmenden kannten die Rolle der Testleitenden, sodass die hohen Werte soziale Erwünschtheit und einen Deckeneffekt enthalten können. Die Anforderungen stützen sich auf eine einzige, auf Empfehlung gewählte Referenzplattform und keine Vorabbefragung. Antwortzeiten und Tragfähigkeit blieben ungeprüft, die Barrierefreiheit ist nur automatisiert und im Ausgangszustand belegt, und getestet wurde ein Prototyp mit 48 vorbereiteten Projekten.

\section{Ausblick}
\label{sec:ausblick}

Die wirkende Einschränkung gehört über die Trefferliste (P17), Handlungen hinter Symbolen und Statusmarken brauchen beschriftete Schaltflächen (P1, P2, P6, P7), und Begriffe sind zu vereinheitlichen. Zu beheben sind außerdem die Kontraste über Bildern und an Feldrändern, die Zielgröße der Personenbilder, der Datenverlust beim Verlassen eines Formulars (P28) und der Schrägstrichbefehl (P30). Offen sind die Suche über Mitwirkende, der Studiengang auf der Profilseite, eine Übersicht der Studiengänge und Module, Untertitel und Transkripte sowie der Abgleich der Löschbefugnisse mit FA-42. Danach wäre nach Nielsen (2000) in mehreren kleinen Runden erneut zu testen, mit mindestens drei Personen je Gruppe und Lehrenden in der Redaktion.

Zwei Anregungen kamen aus der Hochschule selbst. Der Erstprüfer regte an, auf der Startseite zufällig gewählte Projekte aus jeder Fakultät zu zeigen (A. Teufel, persönliche Mitteilung, 8. September 2026), und eine Testperson schlug eine Top 10 nach Fakultäten vor. Ein Portal, das nach Lorenzo und Ittelson (2005) eine Institution repräsentiert, sollte deren Breite schon auf der ersten Seite zeigen; die vorliegende Arbeit schlägt eine kuratierte Top 10 je Fakultät vor, aus der oben wechselnd ein Projekt erscheint. Zweitens schätzten zwei Mitarbeitende der Zentralen Studienberatung in einem Gespräch außerhalb der Erhebung ein, dass Schülerinnen und Schüler vor der Studienwahl die Projekte nicht ohne Weiteres verstehen. Eine Ansicht in verständlicher Sprache, etwa mit Hilfe eines Sprachmodells erzeugt, wäre ein Weg dorthin; zu klären wäre ihr Verhältnis zum Freigabeweg, und weil Gebrauchstauglichkeit kontextgebunden ist (Brooke, 1996), verlangt diese Gruppe einen eigenen Test.

Am Inhaltsmodell wäre die gleichzeitige Bearbeitung durch mehrere Mitwirkende der nächste Schritt, und Kommentare an einzelnen Bausteinen setzen stabile Bausteinkennungen voraus. Für einen Produktivbetrieb fehlen eine serverseitig wirksame Abmeldung, Sperren und Anmeldebegrenzung über mehrere Server sowie die Anonymisierung der Autorschaft beim Löschen eines Kontos. Die Eigenentwicklung trägt zudem nach Lorenzo und Ittelson (2005) das Risiko, das Fachpersonal nicht zu binden; ein Betrieb setzt voraus, dass eine Stelle der Hochschule Wartung und Weiterentwicklung übernimmt, wofür API-Contract, Dokumentation und automatisierte Tests die Grundlage bilden. Offen bleibt schließlich, ob aus der gemessenen Nutzungsabsicht tatsächlich Nutzung wird; das klärt erst eine Erhebung im laufenden Betrieb.
